% !TEX root = ../main.tex
\chapter{PENDAHULUAN}

\section{Latar Belakang}
\label{section:latarbelakang}

% === Paragraf 1: Visi Itera dan Daya Tarik Forest Campus ===
Institut Teknologi Sumatera (Itera) memiliki visi untuk menjadi perguruan
tinggi teknologi yang unggul dengan lingkungan kampus yang nyaman, aman, dan
berwawasan lingkungan \cite{itera2019renstra}. Salah satu bentuk penerapannya
adalah konsep \textit{forest campus}. Konsep ini tidak hanya berkaitan dengan
penataan ruang hijau, tetapi juga dengan keberlanjutan, ketertiban, dan
kualitas hidup sivitas akademika. Dalam pengembangannya, \textit{forest campus}
berhubungan dengan konsep \textit{smart campus}. Konsep ini memanfaatkan sistem
digital, Internet of Things (IoT), dan kecerdasan buatan untuk meningkatkan
efisiensi operasional, pengalaman pengguna, dan keselamatan
\cite{Polin2023SmartCampus}. Perpaduan aspek ekologis dan teknologi tersebut
memperkuat citra Itera sebagai kampus yang mengembangkan lingkungan belajar
yang aman dan berbasis teknologi.

% === Paragraf 2a: Pertumbuhan mahasiswa dan mobilitas sepeda motor ===
Perkembangan kampus tersebut diikuti oleh peningkatan jumlah mahasiswa baru.
Data akademik Itera mencatat tren penerimaan yang meningkat dalam beberapa
tahun terakhir \cite{itera2025dashboard}. Pertumbuhan populasi kampus ini ikut
meningkatkan penggunaan kendaraan pribadi, terutama sepeda motor
\cite{purwanti2018analisisitenas}. Kajian terkait juga menunjukkan peran
mahasiswa dalam inisiatif transportasi dan energi berkelanjutan di perguruan
tinggi \cite{rachmadian2024kesadaran}. Itera telah memulai pengembangan
transportasi umum melalui Smart BRT, tetapi kapasitas dan jangkauannya masih
terbatas \cite{antaranews2025brt}. Selama alternatif transportasi belum
mencakup kebutuhan mobilitas harian, sepeda motor tetap menjadi moda dominan di
pintu gerbang kampus.

Dominasi sepeda motor di lingkungan kampus meningkatkan potensi pelanggaran
lalu lintas, termasuk ketidakpatuhan penggunaan helm. Masalah serupa juga
muncul pada tingkat nasional. Penelitian lintas-seksional pada universitas lain
mencatat 40,8\% mahasiswa pengendara sepeda motor tidak memakai helm secara
konsisten \cite{safitri2020gambaran}. Studi lain menunjukkan bahwa kepatuhan
helm pada remaja dipengaruhi oleh keluarga, kondisi ekonomi, sikap individu,
dan lingkungan sosial-budaya \cite{nobrihas2023kepatuhan}. Untuk mengurangi
risiko kecelakaan, Itera menjalankan program kolaborasi keselamatan berkendara
bersama mitra eksternal. Beberapa perguruan tinggi lain juga menerbitkan
pedoman berlalu lintas secara formal melalui peraturan rektor
\cite{iteranews2024kolaborasi,its2025pedoman}.

Walaupun regulasi dan program keselamatan sudah tersedia, tantangan utama
masih berada pada penegakan di lapangan. Temuan tersebut menunjukkan bahwa
kebijakan tertulis saja belum cukup untuk mengubah perilaku pengendara di
lingkungan kampus. Pengawasan saat ini masih bergantung pada satuan pengamanan
(Satpam) yang bekerja secara manual. Pola ini belum dapat mencakup seluruh
area kampus secara berkelanjutan dan tetap rentan terhadap kesalahan manusia.
Selain itu, pengawasan manual belum menghasilkan data kuantitatif yang
konsisten, sehingga informasi seperti lokasi \textit{hotspot}, waktu puncak,
dan frekuensi pelanggaran belum terdokumentasi dengan baik. Tanpa data
tersebut, Unit Pelaksana Teknis Keselamatan, Kesehatan Kerja, dan Lingkungan
(UPT K3L) Itera sulit menyusun intervensi yang tepat sasaran dan berbasis
bukti. Penelitian ini diarahkan untuk mengisi celah tersebut melalui sistem
otomatis yang merekam pelanggaran helm secara terukur dan mengubahnya menjadi
informasi spasiotemporal yang dapat digunakan oleh pengelola kampus.

% === Paragraf 3: Solusi yang Diusulkan - Sistem Cerdas YOLOv8 + ByteTrack ===
Penelitian ini mengusulkan sistem visi komputer untuk mendeteksi dan
menganalisis pelanggaran penggunaan helm di lingkungan kampus Itera. Rancangan
sistem memadukan YOLOv8 sebagai model deteksi objek dengan ByteTrack sebagai
algoritma pelacakan multiobjek untuk menjaga konsistensi identitas pengendara
antar-\textit{frame}. Kombinasi tersebut telah dievaluasi pada analisis lalu
lintas berbasis video dengan banyak objek bergerak \cite{liu2025vehicle}.
Berbeda dari pendekatan yang berhenti pada identifikasi visual, penelitian ini
mengubah keluaran inferensi menjadi aliran peristiwa (\textit{event stream}).
Aliran tersebut memuat konteks spasial dan temporal, seperti waktu kejadian,
titik kamera, identitas objek, dan status pelanggaran. Setiap peristiwa
dicatat sebagai metadata terstruktur dan disimpan ke \textit{data mart} untuk
mendukung analisis historis, nonvolatil, dan berorientasi analitik. Dengan
rancangan tersebut, sistem tidak hanya berfungsi sebagai alat deteksi otomatis,
tetapi juga sebagai dasar analitik untuk mengidentifikasi \textit{hotspot}
pelanggaran dan menentukan waktu rawan. Hasil akhirnya diharapkan mendukung
pengambilan keputusan keselamatan berbasis data dalam kerangka \textit{smart
campus}.

\section{Rumusan Masalah}
\label{section:rumusanmasalah}

Berdasarkan latar belakang tersebut, rumusan masalah penelitian ini adalah
sebagai berikut.

\begin{enumerate}
  \item Bagaimana Mengembangkan, melatih, dan mengevaluasi model YOLOv8 untuk mendeteksi pengendara sepeda motor berhelm dan tanpa helm pada video CCTV kampus dengan akurasi yang terukur?

  \item Bagaimana mengintegrasikan YOLOv8 dan ByteTrack pada pemrosesan tepi (\textit{edge processing}) serta mendistribusikan keluarannya melalui Apache Kafka dan Apache Spark Structured Streaming menjadi \textit{event stream} pelanggaran helm yang tervalidasi secara spasial (\textit{Intersection over Area}) dan temporal (\textit{N-Frame})?

  \item Bagaimana merancang skema \textit{data mart} berbasis \textit{star schema} dan \textit{dashboard} analitik yang mengubah \textit{event stream} pelanggaran helm menjadi komponen Sistem Pendukung Keputusan (\textit{Decision Support System}) untuk mengidentifikasi pola spasiotemporal pelanggaran di Itera?
\end{enumerate}

\section{Tujuan Penelitian}
\label{section:tujuanpenelitian}

Tujuan penelitian dirumuskan sejajar dengan rumusan masalah sebagai berikut.

\begin{enumerate}
  \item Mengembangkan, melatih, dan mengevaluasi YOLOv8 untuk mendeteksi pengendara sepeda motor berhelm dan tanpa helm pada lingkungan kampus dengan kinerja terukur melalui mAP@50, \textit{precision}, dan \textit{recall}.

  \item Mengimplementasikan integrasi YOLOv8 dan ByteTrack pada pemrosesan tepi serta distribusi data melalui Apache Kafka dan Apache Spark Structured Streaming, dilengkapi filter \textit{Intersection over Area} dan konfirmasi temporal \textit{N-Frame}, untuk menghasilkan \textit{event stream} pelanggaran helm yang tervalidasi.

  \item Merancang dan mengimplementasikan \textit{data mart} berbasis \textit{star schema} pada PostgreSQL beserta \textit{dashboard} analitik berbasis Grafana untuk mengubah \textit{event stream} menjadi Sistem Pendukung Keputusan yang menggambarkan pola spasiotemporal pelanggaran helm di Itera.
\end{enumerate}

\section{Batasan Penelitian}
\label{section:batasanpenelitian}

Ruang lingkup penelitian dibatasi sebagai berikut agar pelaksanaan tetap fokus
dan selesai pada kurun waktu yang ditetapkan.

\begin{enumerate}
  \item Penelitian difokuskan pada deteksi dan analisis pelanggaran penggunaan helm oleh pengendara sepeda motor. Pelanggaran lain, seperti kecepatan, parkir liar, atau rambu, tidak dibahas dalam penelitian ini.

  \item Penelitian berfokus pada analisis historis dan agregatif dari aliran peristiwa pelanggaran helm. Sistem tidak dirancang untuk memprediksi kejadian, memberikan rekomendasi kebijakan otomatis, maupun membandingkan arsitektur \textit{streaming} dengan non-\textit{streaming}.

  \item YOLOv8 digunakan sebagai komponen untuk mengekstraksi informasi visual dari data video. Penelitian tidak membahas modifikasi arsitektur internal YOLOv8, perbandingan dengan detektor lain, maupun optimasi pada tingkat algoritmik.

  \item Pengumpulan dan pengujian data dilakukan dalam rentang waktu tertentu, yaitu pada 5 Februari 2026 dengan jendela observasi pukul 07.00--17.00 WIB.

  \item Evaluasi operasional menggunakan sumber komputasi yang tersedia pada satu mesin pengujian.
\end{enumerate}

% \section{Manfaat Penelitian}
% \label{section:manfaatpenelitian}

% Penelitian ini diharapkan memberikan tiga lapis manfaat yang saling melengkapi.

% \begin{enumerate}
%   \item \textbf{Manfaat akademik.} Studi ini melengkapi literatur \textit{intelligent transportation system} di lingkungan kampus dengan kerangka kerja terintegrasi yang memadukan YOLOv8, ByteTrack, filter \textit{Intersection over Area}, konfirmasi temporal \textit{N-Frame}, dan jalur \textit{stream processing} Kafka--Spark. Hasil eksperimennya dapat menjadi rujukan komparasi bagi penelitian sejenis di Indonesia.

%   \item \textbf{Manfaat praktis.} Sistem yang dibangun siap dimanfaatkan UPT K3L Itera sebagai alat bantu pemantauan harian. \textit{Dashboard} Grafana menyajikan distribusi waktu dan lokasi pelanggaran sehingga unit pengamanan dapat menempatkan personel patroli pada jam dan titik berisiko tinggi tanpa harus menyebar rata sepanjang hari.

%   \item \textbf{Manfaat kebijakan.} Catatan pelanggaran yang persisten di \textit{data mart} menyediakan basis bukti untuk merumuskan Standar Operasional Prosedur (SOP) keselamatan lalu lintas kampus, mengevaluasi efektivitas program sosialisasi helm, dan mendukung pelaporan periodik ke pimpinan institusi.
% \end{enumerate}

% Sub bab lain dapat ditambahkan, misalnya:
%\section{Hipotesis}
